% Generated draft: requires mathematical and source review.
\documentclass[11pt]{article}
\usepackage[margin=25mm]{geometry}
\usepackage{fontspec}
\setmainfont{Latin Modern Roman}
\usepackage{amsmath,amssymb,mathtools,unicode-math}
\setmathfont{Latin Modern Math}
\usepackage{enumitem,needspace,xurl,hyperref}
\setlength{\parindent}{0pt}
\setlength{\parskip}{5pt}
\setlength{\emergencystretch}{3em}
\newcommand{\N}{\mathbb{N}}
\newcommand{\Z}{\mathbb{Z}}
\newcommand{\Q}{\mathbb{Q}}
\newcommand{\R}{\mathbb{R}}
\newcommand{\C}{\mathbb{C}}
\newcommand{\F}{\mathbb{F}}
\newcommand{\A}{\mathbb{A}}
\newcommand{\PP}{\mathbb{P}}
\newcommand{\E}{\mathbb{E}}
\newcommand{\eps}{\varepsilon}
\newcommand{\dd}{\,\mathrm{d}}
\begin{document}
\section*{UnsolvedMath ID: 10000045}
\textbf{Source problem code:} catalogue entry
\subsection{Definitions and mathematical statement}
Consider one-dimensional long-range percolation on \(\Z\), with long-edge probabilities behaving like \(\beta|i-j|^{-2}\). The excerpt says proportional to this expression but does not give an exact probability rule or treatment of short edges. Common alternatives include \(1-\exp(-\beta|i-j|^{-2})\) or a truncated multiple, possibly with nearest-neighbor edges forced present; these are not silently identified.

Write \(D_\beta(0,n)\) for graph distance under the intended connectedness and path conventions. The source describes a deterministic distance exponent \(\theta(\beta)\) by typical distances of order \(n^{\theta(\beta)}\). One precise possible interpretation is
\[
 \frac{\log D_\beta(0,n)}{\log n}\longrightarrow\theta(\beta)
 \quad\text{in probability as }n\to\infty,
\]
assuming finite distances. The excerpt does not specify this mode of convergence, expectations or quantiles, or whether paths are allowed outside \([0,n]\).

The questions are whether \(\theta\) is continuous in \(\beta\), whether it is monotone, and whether a nontrivial metric scaling limit exists. For an edge law increasing in \(\beta\), the natural monotonicity direction is nonincreasing for distances and any consistently defined exponent; whether strict monotonicity is intended must be checked. A scaling-limit statement needs a spatial domain, distance normalization, rooted or unrooted objects, convergence topology and nondegeneracy criterion, none supplied.

No parameter threshold or explicit exponent from the research summary is adopted: that summary revises its own claims and mixes logarithmic, polylogarithmic and polynomial distance regimes without a consistent exact model. The critical inverse-square edge-decay question is preserved separately from those unverified numerical assertions.
\subsection{Short English statement}
For inverse-square long-range percolation on the integer line, determine continuity and monotonicity of the distance exponent and existence of a nondegenerate metric scaling limit. Exact edge laws, distance conventions, and scaling topology need clarification.
\subsection{Sources}
\begin{itemize}
\item Local UnsolvedMath registry, permanent ID 10000045, source code catalogue entry; Itai Benjamini; no proposal year supplied; updated 2026-08-06.
\item Benjamini, \emph{Coarse Geometry and Randomness}, Saint-Flour notes, Open problem 11.21, PDF page 92: \url{https://www.wisdom.weizmann.ac.il/~itai/stflouraug24.pdf#page=92}. Supplied archive: \url{https://arquivo.pt/noFrame/replay/20201231041548id_/http://www.wisdom.weizmann.ac.il/~itai/stflouraug24.pdf}. Registry access date 2026-07-29; extraction checked against rendered PDF, as reported in provenance. Item retained after deduplication and source-side review; bounded audit found no complete resolution. Current status and reuse terms require review.
\item The supplied summary names Benjamini--Berger, Berger, Ding, Biskup and Hutchcroft, and approximate 2020--2023 work, but supplies no precise papers or consistent quantitative theorem. Its status claims were not independently checked.
\end{itemize}
\end{document}
